\begin{table}[htbp]
  \centering
  \caption{The comparability protocol, stage by stage. Each row pairs a generic benchmarking principle with its concrete manifestation for relational database access layers, the failure that follows if the stage is omitted, and the evidence that the stage was load-bearing in this study. We do not claim to originate these principles; in the sources searched we found no prior access-layer benchmark applying all five together. Throughputs abbreviated PG (PostgreSQL) and MySQL; ``conns'' denotes client connections.}
  \label{tab:protocol_mapping}
  \begin{adjustbox}{max width=\textwidth}
  \footnotesize
  \begin{tabular}{p{1.7cm} p{2.3cm} p{2.5cm} p{3.3cm} p{2.5cm} p{3.4cm}}
    \toprule
    \textbf{Stage} & \textbf{Benchmark decision} & \textbf{Generic principle} & \textbf{Access-layer manifestation} & \textbf{Failure if omitted} & \textbf{Evidence it was load-bearing} \\
    \midrule
    Correctness oracle &
    Which treatments are timed at all &
    Validity precedes speed; never time incorrect results &
    Layers differ in output shape and mutation semantics; a ``fast'' cell may be a silent error doing less work, whereas engines return one canonical result &
    Broken/incomplete layers look fastest and top the ranking &
    Broken MikroORM MySQL write led the insert ranking until the write-state gate excluded it (silent 500s); a read cross-check caught a fan-out aggregation bug and a driver result-shape mismatch \\
    \addlinespace
    Treatment-definition rule &
    Which API and loading strategy of each layer is measured &
    Pre-register and fix the system-under-test configuration reproducibly &
    One library exposes many equally-official query paths with different round-trip counts; a pure engine has one query language, so no selection rule is needed &
    Arbitrary/cherry-picked API choice; spread reflects the tester, not the defaults &
    Documentation-selected deep fetch spreads $7.15\times$ (PG) / $4.96\times$ (MySQL); round-trips differ 1--4 by design; Drizzle is the documented tie-break case \\
    \addlinespace
    Strategy standardization &
    Add a standardized same-SQL contrast (a compound diagnostic, not a mechanism decomposition) &
    Hold confounds constant via a common-baseline contrast &
    Access-layer time mixes loading strategy with execution machinery; engines have no ORM loading strategy to confound, so no same-SQL baseline is meaningful &
    Mis-read the whole documentation-selected spread as intrinsic library ``overhead'' &
    The $7.15\times$/$4.96\times$ spread narrows to $1.68\times$ (PG) / $2.01\times$ (MySQL) among raw paths on common SQL, so most of the documentation-selected spread disappears under this compound raw-SQL standardization; which component is responsible is not identified \\
    \addlinespace
    Capacity identification &
    Locate each layer's saturating throughput before fixing an operating point &
    Interpret latency relative to capacity, not at an arbitrary fixed load &
    Layers saturate at very different throughputs, so a fixed connection count lands them at different capacity fractions; a single engine saturates near one throughput &
    A fixed load compares layers at unequal utilization &
    A per-pattern sweep (1--200) puts every knee at or below 50 conns; median utilization 97--100\% at the 50-conn point; native $\sim$$7\times$ the slowest ORM's capacity \\
    \addlinespace
    Demand \& utilization &
    Report tail at equal demand and at matched own-capacity fractions, as distinct quantities &
    Distinguish demand from utilization; correct coordinated omission &
    Because layer capacities differ, an equal-demand tail conflates queueing with intrinsic latency; engines at similar capacity make the two nearly coincide &
    Read a capacity/queueing gap as an intrinsic latency difference &
    At 50 conns deep-fetch p99 runs pg 20\,ms to MikroORM 116\,ms; at 50\% of each layer's own capacity the tails converge to $\sim$2--5\,ms, so the gap is capacity-and-queueing, not intrinsic latency \\
    \bottomrule
  \end{tabular}
  \end{adjustbox}
\end{table}